\section{Synthetic Respondents as Measurement Instruments}\label{sec:background}

\paragraph{Synthetic respondents and the criterion they need.} A language model conditioned on a description of a population can be asked to stand in for a sample of its members. In silicon sampling the description enters the prompt and the model's answer distribution is read as an estimate of the population's \citep{argyle2023}; in the homo silicus program the model plays an economic agent in a stylized decision problem \citep{horton2023}; a systematic review maps how fast these uses have spread across survey research \citep{buskirk2026}. The evidence on prompted models is summarized in the introduction. Fine-tuning on survey responses is one response to it \citep{suh2025,cao2025,krsteski2026}. Two reviews in this journal draw the methodological conclusion that matters here. \citet{kozlowski2025} conclude that validation against human data remains essential, and \citet{broska2025} propose a design in which model predictions are potentially informative observations and human subjects remain the gold standard. Both require a criterion against which the synthetic reading can fail. Our evaluation supplies one criterion for each link of the chain.

\paragraph{Construct, measure, and anchor.} Three objects have to be kept apart. A construct is a theoretical preference attribute, such as generalized trust. A measure operationalizes it through observed answers, such as the response distribution to a WVS question about whether most people can be trusted. An anchor is an aggregate score that conditions the instrument, such as a country's GPS trust coordinate. The GPS supplies anchors for six economic preferences (trust, risk taking, patience, altruism, positive reciprocity and negative reciprocity), measured with experimentally validated modules in a common design across countries \citep{falk2018}. We treat these coordinates as declared aggregate moments: the researcher chooses them, and nothing in the construction certifies them as the right moments for a given survey. Mapping a WVS item to a GPS dimension makes the item a candidate indicator, and an average of such items is a proposed composite until its interpretation is supported. Agreement between instrument and anchor, between instrument and people, and between anchor and people across countries are therefore three different facts, and they are the three comparisons that the chain makes. The multitrait--multimethod logic of \citet{campbellfiske1959} applies directly: the human survey is the other method, and an instrument that agrees with its own anchor has passed a construction check without yet meeting a criterion.

\paragraph{Cross-national comparability.} The unit of analysis is the country. Country-level agreement does not establish that an item measures the same construct for individuals in different countries. Loadings and intercepts can differ across groups in ways that leave country means comparable for some purposes and misleading for others \citep{steenkamp1998,davidov2014}. Work in this journal argues that composite scores aggregated to the country level can suffice for simple group comparisons despite noninvariance, while understating uncertainty \citep{bratt2025}, and that unaddressed measurement error and, less so, noninvariance bias multilevel estimates \citep{pokropek2025}. We adopt the same posture. Every comparison with respondents below is a group-level directional question about country rankings, with no claim about invariance or individual-level validity.

\paragraph{The empty cell.} Prior evaluations score models on questions whose human answers are known, usually with the population named in the prompt, and fine-tuning approaches learn from human survey responses. The question here is, link by link, what an instrument built from declared aggregate inputs alone preserves of its anchor and of the respondents' answers, with the selection rule and inference for new items registered before the human criterion is read and with comparators that read only the questionnaire. Section~\ref{sec:design} describes such an evaluation and Section~\ref{sec:results} reports it.
